% --- IMPORTS ---
\documentclass[leqno]{article}
\usepackage{graphicx}
\usepackage{dsfont}
\usepackage{mathtools}
\usepackage{amssymb}
\usepackage{amsmath}
\usepackage{amsthm}
\usepackage{float}
\usepackage{ragged2e}
\usepackage{tensor}
\usepackage{fancyhdr}
\usepackage{empheq}
\usepackage[most]{tcolorbox}
\usepackage{enumitem}
\usepackage{dirtytalk}
\usepackage{marginnote}
\usepackage{microtype}
\usepackage[
  a4paper,
  left=0.8in,
  right=1in,
  top=1in,
  bottom=0.5in,
  headheight=14pt,
  headsep=0.4in
]{geometry}
\setlength{\parindent}{0cm}
% --- TikZ Packages ---
\usepackage{pgfplots}
\pgfplotsset{compat=1.18}
\usepgfplotslibrary{fillbetween, polar}
\usetikzlibrary{calc, positioning}
\usetikzlibrary{angles, quotes}
\usetikzlibrary{arrows.meta}
\usetikzlibrary{decorations.pathreplacing, decorations.markings}
\usetikzlibrary{patterns}
\usetikzlibrary{intersections}
% --- URL Packages ---
\usepackage[colorlinks=true,linkcolor=red,citecolor=magenta,urlcolor=black]{hyperref}%
\urlstyle{same}
% --- END OF IMPORTS ---

% --- CUSTOM COMMANDS ---
\RenewDocumentCommand{\marks}{o m}{%
  \IfValueTF{#1}%
    {\marginnote{\raggedright\textbf{[#2]}}[#1]}%
    {\marginnote{\raggedright\textbf{[#2]}}}%
}
\newcommand{\vspacerule}[2]{%
  \par\vspace*{#1}%
  \noindent\hrulefill\par
  \vspace*{#2}%
}
\definecolor{scarlet}{RGB}{205, 40, 75}
\newcommand{\questionitem}{%
  \item\phantomsection
  \addcontentsline{toc}{section}{Question \number\value{enumi}}%
}
\renewcommand{\contentsname}{Questions}
% --- END OF CUSTOM COMMANDS ---

% --- HEADER CONFIG ---
\pagestyle{fancy}
\fancyhead{}
\fancyfoot{}
\renewcommand{\headrulewidth}{0.9pt}
\lhead{\textit{Solutions} - Proof by Induction - Derivatives}
\chead{}
\rhead{\href{https://danieljsmith.org}{danieljsmith.org}}
% --- END OF HEADER CONFIG ---

\begin{document}
\vspace*{20pt}
\tableofcontents
\newpage
\begin{enumerate}[
  label=\textbf{\arabic*.},
  leftmargin=*,
  topsep=0pt,
  partopsep=0pt,
  parsep=0pt
]

% ---- Question 1 ----
\questionitem

It is given that \[y = x^2 e^{ax}\] where $a$ is a constant. \\

Prove by mathematical induction that, for all integers $n \ge 2$,
\[
\frac{\mathrm{d}^n y}{\mathrm{d}x^n} = a^{n-2}\left(a^2x^2 + 2nax + n(n-1)\right)e^{ax}
\] \marks[-30pt]{6}

\vspace*{10pt}

\begin{tcolorbox}[colback=white, colframe=scarlet, title={\textbf{Solution}}, breakable, grow to left by=6mm, grow to right by=10mm]
Let
\[
P(n): \quad \frac{\mathrm{d}^n y}{\mathrm{d}x^n}=a^{n-2}\left(a^2x^2+2nax+n(n-1)\right)e^{ax}
\]
We prove \(P(n)\) is true for all integers \(n \ge 2\) by mathematical induction. \\

\textbf{Base case:} \(n=2\)

Since \(y=x^2e^{ax}\),
\begin{align*}
\frac{\mathrm{d}y}{\mathrm{d}x}
&=2xe^{ax}+x^2(ae^{ax}) \\
&=(ax^2+2x)e^{ax}
\end{align*}

Differentiate again:
\begin{align*}
\frac{\mathrm{d}^2y}{\mathrm{d}x^2}
&=\frac{\mathrm{d}}{\mathrm{d}x}\left((ax^2+2x)e^{ax}\right) \\
&=(2ax+2)e^{ax}+(ax^2+2x)(ae^{ax}) \\
&=(2ax+2+a^2x^2+2ax)e^{ax} \\
&=(a^2x^2+4ax+2)e^{ax}
\end{align*}

Now substitute \(n=2\) into the given formula:
\begin{align*}
a^{2-2}\left(a^2x^2+2(2)ax+2(2-1)\right)e^{ax}
&=a^0(a^2x^2+4ax+2)e^{ax} \\
&=(a^2x^2+4ax+2)e^{ax}
\end{align*}

This matches \(\dfrac{\mathrm{d}^2y}{\mathrm{d}x^2}\), so \(P(2)\) is true. \\

\textbf{Inductive Step}\\
Now assume \(P(k)\) is true for some integer \(k\ge2\).  
That is, assume
\[
\frac{\mathrm{d}^k y}{\mathrm{d}x^k}
=a^{k-2}\left(a^2x^2+2kax+k(k-1)\right)e^{ax}
\]
This is the induction hypothesis.

We must show that \(P(k+1)\) is true.

Differentiate the induction hypothesis:
\begin{align*}
\frac{\mathrm{d}^{k+1}y}{\mathrm{d}x^{k+1}}
&=a^{k-2}\frac{\mathrm{d}}{\mathrm{d}x}\left[\left(a^2x^2+2kax+k(k-1)\right)e^{ax}\right]
\end{align*}

Using the product rule,
\begin{align*}
\frac{\mathrm{d}^{k+1}y}{\mathrm{d}x^{k+1}}
&=a^{k-2}\left[\left(2a^2x+2ka\right)e^{ax}+\left(a^2x^2+2kax+k(k-1)\right)ae^{ax}\right] \\
&=a^{k-2}\left[2a^2x+2ka+a^3x^2+2ka^2x+ak(k-1)\right]e^{ax}
\end{align*}

Now simplify the bracket:
\begin{align*}
2a^2x+2ka+a^3x^2+2ka^2x+ak(k-1)
&=a^3x^2+(2+2k)a^2x+\bigl(2k+k(k-1)\bigr)a \\
&=a^3x^2+2(k+1)a^2x+k(k+1)a
\end{align*}

So
\begin{align*}
\frac{\mathrm{d}^{k+1}y}{\mathrm{d}x^{k+1}}
&=a^{k-2}\left(a^3x^2+2(k+1)a^2x+k(k+1)a\right)e^{ax} \\
&=a^{k-1}\left(a^2x^2+2(k+1)ax+k(k+1)\right)e^{ax}
\end{align*}

Since
\[
k(k+1)=(k+1)\bigl((k+1)-1\bigr)
\]
this becomes
\[
\frac{\mathrm{d}^{k+1}y}{\mathrm{d}x^{k+1}}
=a^{(k+1)-2}\left(a^2x^2+2(k+1)ax+(k+1)\bigl((k+1)-1\bigr)\right)e^{ax}
\]

This is exactly \(P(k+1)\), so the inductive step is complete.\\


\textbf{Conclusion:}

Therefore, since \(P(2)\) is true and \(P(k)\Rightarrow P(k+1)\), it follows by mathematical induction that, for all integers \(n\ge2\),
\[
\frac{\mathrm{d}^n y}{\mathrm{d}x^n}=a^{n-2}\left(a^2x^2+2nax+n(n-1)\right)e^{ax}
\]
\end{tcolorbox}


\newpage

% ---- Question 2 ----
\questionitem

Prove by mathematical induction that, for every positive integer $n$,
\[
\frac{\mathrm{d}^n}{\mathrm{d}x^n}(x^3 e^{-x}) = (-1)^n e^{-x}\bigl(x^3 - 3n x^2 + 3n(n-1)x - n(n-1)(n-2)\bigr)
\] \marks[-30pt]{7}

\vspace*{10pt}

\begin{tcolorbox}[colback=white, colframe=scarlet, title={\textbf{Solution}}, breakable, grow to left by=6mm, grow to right by=10mm]
We prove the result by mathematical induction.

Let
\[
P(n):\quad \frac{\mathrm{d}^n}{\mathrm{d}x^n}(x^3e^{-x})
=
(-1)^n e^{-x}\bigl(x^3-3nx^2+3n(n-1)x-n(n-1)(n-2)\bigr)
\]

We must show that $P(n)$ is true for every positive integer $n$.

\textbf{Base case: $n=1$.}

Differentiate $x^3e^{-x}$ using the product rule:
\begin{align*}
\frac{\mathrm{d}}{\mathrm{d}x}(x^3e^{-x})
&=3x^2e^{-x}+x^3(-e^{-x})\\
&=e^{-x}(3x^2-x^3)\\
&=(-1)e^{-x}(x^3-3x^2)
\end{align*}

Now substitute $n=1$ into the required formula:
\[
(-1)^1 e^{-x}\bigl(x^3-3(1)x^2+3(1)(0)x-(1)(0)(-1)\bigr)
=
(-1)e^{-x}(x^3-3x^2)
\]

This matches the derivative found above, so $P(1)$ is true.

\textbf{Inductive hypothesis.}

Assume that $P(k)$ is true for some positive integer $k$.

That is, assume
\[
\frac{\mathrm{d}^k}{\mathrm{d}x^k}(x^3e^{-x})
=
(-1)^k e^{-x}\bigl(x^3-3kx^2+3k(k-1)x-k(k-1)(k-2)\bigr)
\]

For convenience, let
\[
P_k(x)=x^3-3kx^2+3k(k-1)x-k(k-1)(k-2)
\]

Then the inductive hypothesis is
\[
\frac{\mathrm{d}^k}{\mathrm{d}x^k}(x^3e^{-x})=(-1)^k e^{-x}P_k(x)
\]

\textbf{Inductive step.}

We now show that $P(k+1)$ is true.

Differentiate the inductive hypothesis:
\begin{align*}
\frac{\mathrm{d}^{k+1}}{\mathrm{d}x^{k+1}}(x^3e^{-x})
&=
(-1)^k\frac{\mathrm{d}}{\mathrm{d}x}\bigl(e^{-x}P_k(x)\bigr)\\
&=
(-1)^k\left((-e^{-x})P_k(x)+e^{-x}P_k'(x)\right)\\
&=
(-1)^k e^{-x}\bigl(P_k'(x)-P_k(x)\bigr)\\
&=
(-1)^{k+1} e^{-x}\bigl(P_k(x)-P_k'(x)\bigr)
\end{align*}

Now find $P_k'(x)$:
\[
P_k'(x)=3x^2-6kx+3k(k-1)
\]

So
\begin{align*}
P_k(x)-P_k'(x)
&=
\bigl(x^3-3kx^2+3k(k-1)x-k(k-1)(k-2)\bigr)\\
&\qquad-\bigl(3x^2-6kx+3k(k-1)\bigr)\\
&=
x^3-3(k+1)x^2+\bigl(3k(k-1)+6k\bigr)x\\
&\qquad-\bigl(k(k-1)(k-2)+3k(k-1)\bigr)\\
&=
x^3-3(k+1)x^2+3k(k+1)x-k(k-1)(k+1)\\
&=
x^3-3(k+1)x^2+3(k+1)k\,x-(k+1)k\bigl((k+1)-2\bigr)
\end{align*}

Therefore
\[
\frac{\mathrm{d}^{k+1}}{\mathrm{d}x^{k+1}}(x^3e^{-x})
=
(-1)^{k+1}e^{-x}\bigl(x^3-3(k+1)x^2+3(k+1)k\,x-(k+1)k((k+1)-2)\bigr)
\]

This is exactly the required formula with $n=k+1$, so $P(k+1)$ is true.

\textbf{Conclusion.}

Since $P(1)$ is true, and $P(k)\Rightarrow P(k+1)$, the statement holds for all positive integers $n$ by mathematical induction.

Hence
\[
\frac{\mathrm{d}^n}{\mathrm{d}x^n}(x^3 e^{-x}) = (-1)^n e^{-x}\bigl(x^3 - 3n x^2 + 3n(n-1)x - n(n-1)(n-2)\bigr)
\]

for every positive integer $n$.
\end{tcolorbox}


\newpage

% ---- Question 3 ----
\questionitem

Given that
\[
y=e^{-2x}\cosh 3x
\]
prove by induction that for $n \in \mathbb{N}$
\[
\frac{\mathrm{d}^n y}{\mathrm{d}x^n}=e^{-2x}\left(\frac{1+(-5)^n}{2}\cosh 3x+\frac{1-(-5)^n}{2}\sinh 3x\right)
\] \marks[-30pt]{6}

\vspace*{10pt}

\begin{tcolorbox}[colback=white, colframe=scarlet, title={\textbf{Solution}}, breakable, grow to left by=6mm, grow to right by=10mm]
We prove the result by induction on $n$.

Let $P(n)$ be the statement
\[
\frac{\mathrm{d}^n y}{\mathrm{d}x^n}
=
e^{-2x}\left(\frac{1+(-5)^n}{2}\cosh 3x+\frac{1-(-5)^n}{2}\sinh 3x\right)
\]

where
\[
y=e^{-2x}\cosh 3x
\]

\textbf{Base case: $n=1$}

Differentiate $y$ using the product rule. Also,
\[
\frac{\mathrm{d}}{\mathrm{d}x}\big(\cosh 3x\big)=3\sinh 3x
\]

So
\begin{align*}
\frac{\mathrm{d}y}{\mathrm{d}x}
&=\frac{\mathrm{d}}{\mathrm{d}x}\big(e^{-2x}\cosh 3x\big) \\
&=e^{-2x}(-2)\cosh 3x+e^{-2x}(3\sinh 3x) \\
&=e^{-2x}\big(-2\cosh 3x+3\sinh 3x\big)
\end{align*}

Now check this against the required formula when $n=1$:
\[
\frac{1+(-5)^1}{2}=\frac{1-5}{2}=-2,
\qquad
\frac{1-(-5)^1}{2}=\frac{1+5}{2}=3
\]

So
\[
\frac{\mathrm{d}y}{\mathrm{d}x}
=
e^{-2x}\left(\frac{1+(-5)^1}{2}\cosh 3x+\frac{1-(-5)^1}{2}\sinh 3x\right)
\]

Hence $P(1)$ is true.

\textbf{Inductive hypothesis}

Assume that $P(n)$ is true for some $n\in\mathbb{N}$. That is, assume
\[
\frac{\mathrm{d}^n y}{\mathrm{d}x^n}
=
e^{-2x}\left(\frac{1+(-5)^n}{2}\cosh 3x+\frac{1-(-5)^n}{2}\sinh 3x\right)
\]

For simplicity, write
\[
A=\frac{1+(-5)^n}{2},
\qquad
B=\frac{1-(-5)^n}{2}
\]

Then the inductive hypothesis becomes
\[
\frac{\mathrm{d}^n y}{\mathrm{d}x^n}=e^{-2x}(A\cosh 3x+B\sinh 3x)
\]

\textbf{Inductive step}

We now differentiate to find $\dfrac{\mathrm{d}^{n+1}y}{\mathrm{d}x^{n+1}}$:
\begin{align*}
\frac{\mathrm{d}^{n+1}y}{\mathrm{d}x^{n+1}}
&=\frac{\mathrm{d}}{\mathrm{d}x}\left[e^{-2x}(A\cosh 3x+B\sinh 3x)\right] \\
&=-2e^{-2x}(A\cosh 3x+B\sinh 3x)
+e^{-2x}(3A\sinh 3x+3B\cosh 3x) \\
&=e^{-2x}\left[(-2A+3B)\cosh 3x+(3A-2B)\sinh 3x\right]
\end{align*}

Now simplify the coefficients.

First,
\begin{align*}
-2A+3B
&=-2\left(\frac{1+(-5)^n}{2}\right)+3\left(\frac{1-(-5)^n}{2}\right) \\
&=\frac{-2-2(-5)^n+3-3(-5)^n}{2} \\
&=\frac{1-5(-5)^n}{2} \\
&=\frac{1+(-5)^{n+1}}{2}
\end{align*}

since
\[
(-5)^{n+1}=-5(-5)^n
\]

Also,
\begin{align*}
3A-2B
&=3\left(\frac{1+(-5)^n}{2}\right)-2\left(\frac{1-(-5)^n}{2}\right) \\
&=\frac{3+3(-5)^n-2+2(-5)^n}{2} \\
&=\frac{1+5(-5)^n}{2} \\
&=\frac{1-(-5)^{n+1}}{2}
\end{align*}

Therefore
\[
\frac{\mathrm{d}^{n+1}y}{\mathrm{d}x^{n+1}}
=
e^{-2x}\left(\frac{1+(-5)^{n+1}}{2}\cosh 3x+\frac{1-(-5)^{n+1}}{2}\sinh 3x\right)
\]

So $P(n+1)$ is true.

Since $P(1)$ is true, and $P(n)\Rightarrow P(n+1)$, the result follows by induction for all $n\in\mathbb{N}$.

Hence
\[
\frac{\mathrm{d}^n y}{\mathrm{d}x^n}=e^{-2x}\left(\frac{1+(-5)^n}{2}\cosh 3x+\frac{1-(-5)^n}{2}\sinh 3x\right)
\quad \text{for all } n\in\mathbb{N}
\]
\end{tcolorbox}


\newpage

% ---- Question 4 ----
\questionitem

Let \[y = x^2\cosh x\]

Prove by induction that, for all integers $n \ge 1$,
\[
\frac{\mathrm{d}^{2n}y}{\mathrm{d}x^{2n}} = x^2\cosh x + 4nx\sinh x + 2n(2n-1)\cosh x
\] \marks[-20pt]{6}

\vspace*{10pt}

\begin{tcolorbox}[colback=white, colframe=scarlet, title={\textbf{Solution}}, breakable, grow to left by=6mm, grow to right by=10mm]
Let
\[
P(n): \quad \frac{\mathrm{d}^{2n}y}{\mathrm{d}x^{2n}}=x^2\cosh x+4nx\sinh x+2n(2n-1)\cosh x
\]

We prove \(P(n)\) by mathematical induction for all integers \(n\ge 1\).

\textbf{Base case: \(n=1\).}

Since \(y=x^2\cosh x\),
\begin{align*}
\frac{\mathrm{d}y}{\mathrm{d}x}
&=2x\cosh x+x^2\sinh x
\end{align*}
and so
\begin{align*}
\frac{\mathrm{d}^2y}{\mathrm{d}x^2}
&=\frac{\mathrm{d}}{\mathrm{d}x}\bigl(2x\cosh x+x^2\sinh x\bigr)\\
&=(2\cosh x+2x\sinh x)+(2x\sinh x+x^2\cosh x)\\
&=x^2\cosh x+4x\sinh x+2\cosh x
\end{align*}

This is
\[
x^2\cosh x+4(1)x\sinh x+2\cdot 1(2\cdot 1-1)\cosh x
\]
so \(P(1)\) is true.

\textbf{Inductive step.}

Assume \(P(k)\) is true for some integer \(k\ge 1\). Then
\[
\frac{\mathrm{d}^{2k}y}{\mathrm{d}x^{2k}}=x^2\cosh x+4kx\sinh x+2k(2k-1)\cosh x
\]

We must show that \(P(k+1)\) is true.

Differentiate once:
\begin{align*}
\frac{\mathrm{d}^{2k+1}y}{\mathrm{d}x^{2k+1}}
&=\frac{\mathrm{d}}{\mathrm{d}x}\bigl(x^2\cosh x+4kx\sinh x+2k(2k-1)\cosh x\bigr)\\
&=(2x\cosh x+x^2\sinh x)+4k(\sinh x+x\cosh x)+2k(2k-1)\sinh x\\
&=x^2\sinh x+(4k+2)x\cosh x+\bigl(4k+2k(2k-1)\bigr)\sinh x\\
&=x^2\sinh x+(4k+2)x\cosh x+2k(2k+1)\sinh x
\end{align*}

Differentiate again:
\begin{align*}
\frac{\mathrm{d}^{2k+2}y}{\mathrm{d}x^{2k+2}}
&=\frac{\mathrm{d}}{\mathrm{d}x}\bigl(x^2\sinh x+(4k+2)x\cosh x+2k(2k+1)\sinh x\bigr)\\
&=(2x\sinh x+x^2\cosh x)+(4k+2)(\cosh x+x\sinh x)+2k(2k+1)\cosh x\\
&=x^2\cosh x+\bigl(2+(4k+2)\bigr)x\sinh x+\bigl((4k+2)+2k(2k+1)\bigr)\cosh x\\
&=x^2\cosh x+4(k+1)x\sinh x+\bigl(4k+2+4k^2+2k\bigr)\cosh x\\
&=x^2\cosh x+4(k+1)x\sinh x+(4k^2+6k+2)\cosh x
\end{align*}

Now
\[
4k^2+6k+2=2(k+1)(2k+1)=2(k+1)\bigl(2(k+1)-1\bigr)
\]

So
\begin{align*}
\frac{\mathrm{d}^{2k+2}y}{\mathrm{d}x^{2k+2}}
&=x^2\cosh x+4(k+1)x\sinh x+2(k+1)\bigl(2(k+1)-1\bigr)\cosh x
\end{align*}

This is exactly \(P(k+1)\). Therefore \(P(k)\Rightarrow P(k+1)\).

Since \(P(1)\) is true, and \(P(k)\Rightarrow P(k+1)\), it follows by induction that for all integers \(n\ge 1\),
\[
\frac{\mathrm{d}^{2n}y}{\mathrm{d}x^{2n}}=x^2\cosh x+4nx\sinh x+2n(2n-1)\cosh x
\]
\end{tcolorbox}


\newpage

% ---- Question 5 ----
\questionitem

A function is defined by $y=f(t)$ where $f(t)=(1+at)\ln(1+at)$ and $a$ is a constant. \\

\begin{enumerate}[label=\textbf{(\alph*)}]
  \item By considering $\frac{\mathrm{d}^2y}{\mathrm{d}t^2}$, $\frac{\mathrm{d}^3y}{\mathrm{d}t^3}$, $\frac{\mathrm{d}^4y}{\mathrm{d}t^4}$ and $\frac{\mathrm{d}^5y}{\mathrm{d}t^5}$, make a conjecture for a general formula for $\frac{\mathrm{d}^n y}{\mathrm{d}t^n}$ in terms of $n$, $a$ and $t$ for any integer $n \ge 2$. \marks{3} \\
  \item Use induction to prove the formula conjectured in part $(a)$. \marks{4} \\
  \item In the case when $f(t)=(1+3t)\ln(1+3t)$, find the rate at which the $6^{\mathrm{th}}$ derivative of $f(t)$ is varying when $t=\frac{5}{3}$. \marks{2} \\
\end{enumerate}

\vspace*{10pt}

\begin{tcolorbox}[colback=white, colframe=scarlet, title={\textbf{Solution}}, breakable, grow to left by=6mm, grow to right by=10mm]
\begin{enumerate}[label=\textbf{(\alph*)}]
  \item Let
  \[
  y=(1+at)\ln(1+at)
  \]

  First differentiate once:
  \begin{align*}
  \frac{\mathrm{d}y}{\mathrm{d}t}
  &= \frac{\mathrm{d}}{\mathrm{d}t}(1+at)\cdot \ln(1+at) + (1+at)\cdot \frac{\mathrm{d}}{\mathrm{d}t}\bigl(\ln(1+at)\bigr) \\
  &= a\ln(1+at) + (1+at)\cdot \frac{a}{1+at} \\
  &= a\ln(1+at)+a
  \end{align*}

  Now find the next derivatives:
  \begin{align*}
  \frac{\mathrm{d}^2y}{\mathrm{d}t^2}
  &= \frac{\mathrm{d}}{\mathrm{d}t}\bigl(a\ln(1+at)+a\bigr)
  = a\cdot \frac{a}{1+at}
  = \frac{a^2}{1+at}
  \\
  \frac{\mathrm{d}^3y}{\mathrm{d}t^3}
  &= \frac{\mathrm{d}}{\mathrm{d}t}\left(a^2(1+at)^{-1}\right)
  = a^2(-1)(1+at)^{-2}\cdot a
  = -\frac{a^3}{(1+at)^2}
  \\
  \frac{\mathrm{d}^4y}{\mathrm{d}t^4}
  &= \frac{\mathrm{d}}{\mathrm{d}t}\left(-a^3(1+at)^{-2}\right)
  = -a^3(-2)(1+at)^{-3}\cdot a
  = \frac{2a^4}{(1+at)^3}
  \\
  \frac{\mathrm{d}^5y}{\mathrm{d}t^5}
  &= \frac{\mathrm{d}}{\mathrm{d}t}\left(2a^4(1+at)^{-3}\right)
  = 2a^4(-3)(1+at)^{-4}\cdot a
  = -\frac{6a^5}{(1+at)^4}
  \end{align*}

  The pattern is:
  \[
  \frac{\mathrm{d}^2y}{\mathrm{d}t^2}=\frac{a^2}{(1+at)^1},\quad
  \frac{\mathrm{d}^3y}{\mathrm{d}t^3}=-\frac{a^3}{(1+at)^2},\quad
  \frac{\mathrm{d}^4y}{\mathrm{d}t^4}=\frac{2a^4}{(1+at)^3},\quad
  \frac{\mathrm{d}^5y}{\mathrm{d}t^5}=-\frac{6a^5}{(1+at)^4}
  \]

  and the numbers \(1,1,2,6\) are \(0!,1!,2!,3!\).

  So the conjecture is
  \[
  f^{(n)}(t)=\frac{(-1)^n(n-2)!a^n}{(1+at)^{\,n-1}}
  \qquad \text{for all integers } n\ge 2
  \]

  \item We prove by induction that
  \[
  f^{(n)}(t)=\frac{(-1)^n(n-2)!a^n}{(1+at)^{\,n-1}}
  \qquad (n\ge 2)
  \]

  \textbf{Base case: \(n=2\).}

  From part (a),
  \[
  f''(t)=\frac{a^2}{1+at}
  \]

  The formula gives
  \[
  \frac{(-1)^2(2-2)!a^2}{(1+at)^{2-1}}
  =\frac{1\cdot 0!\,a^2}{1+at}
  =\frac{a^2}{1+at}
  \]

  so the result is true for \(n=2\).

  \textbf{Induction hypothesis.}

  Assume that for some integer \(k\ge 2\),
  \[
  f^{(k)}(t)=\frac{(-1)^k(k-2)!a^k}{(1+at)^{\,k-1}}
  \]

  \textbf{Induction step.}

  Differentiate both sides:
  \begin{align*}
  f^{(k+1)}(t)
  &= \frac{\mathrm{d}}{\mathrm{d}t}\left[(-1)^k(k-2)!a^k(1+at)^{-(k-1)}\right] \\
  &= (-1)^k(k-2)!a^k \cdot \left(-(k-1)(1+at)^{-k}\cdot a\right) \\
  &= (-1)^{k+1}(k-1)(k-2)!a^{k+1}(1+at)^{-k} \\
  &= \frac{(-1)^{k+1}(k-1)!a^{k+1}}{(1+at)^k}
  \end{align*}

  This is exactly the required formula with \(n=k+1\), since
  \[
  (k+1)-2 = k-1
  \qquad \text{and} \qquad
  (k+1)-1 = k
  \]

  So, if the formula is true for \(n=k\), it is also true for \(n=k+1\).

  Since it is true for \(n=2\), it follows by induction that
  \[
  f^{(n)}(t)=\frac{(-1)^n(n-2)!a^n}{(1+at)^{\,n-1}}
  \qquad \text{for all integers } n\ge 2
  \]

  \item Here \(a=3\), so
  \[
  f^{(n)}(t)=\frac{(-1)^n(n-2)!3^n}{(1+3t)^{\,n-1}}
  \]

  The rate at which the \(6^\text{th}\) derivative is varying is the \(7^\text{th}\) derivative:
  \[
  f^{(7)}(t)=\frac{(-1)^7(7-2)!3^7}{(1+3t)^6}
  =-\frac{5!\,3^7}{(1+3t)^6}
  \]

  When \(t=\frac{5}{3}\),
  \[
  1+3t=1+3\cdot \frac{5}{3}=6
  \]

  Hence
  \begin{align*}
  f^{(7)}\left(\frac{5}{3}\right)
  &= -\frac{120\cdot 3^7}{6^6} \\
  &= -\frac{120\cdot 3^7}{(2\cdot 3)^6} \\
  &= -\frac{120\cdot 3}{2^6} \\
  &= -\frac{360}{64}
  = -\frac{45}{8}
  \end{align*}

  Therefore, the rate is
  \[
  -\frac{45}{8}
  \]
\end{enumerate}
\end{tcolorbox}


\newpage

% ---- Question 6 ----
\questionitem

The function $f(x)$ is defined by $f(x)=(1+x)^{-\frac12}$, for $x>-1$. \\

\begin{enumerate}[label=\textbf{(\alph*)}]
  \item Prove by mathematical induction that the $n$th derivative of $f(x)$, $f^{(n)}(x)$, for all $n \ge 1$, is given by
  \[
  f^{(n)}(x)=\frac{(-1)^n(2n)!}{2^{2n}n!(1+x)^{n+\frac12}}
  \]
  \marks[-20pt]{4}
  \item Hence prove that
  \[
  (1+x)^{-\frac12}=1-\frac{x}{2}+\frac{3x^2}{8}-\frac{5x^3}{16}+\ldots+(-1)^n\frac{(2n)!}{2^{2n}(n!)^2}x^n+\ldots
  \]
  for $-1<x\le 1$. \marks{3} \\
\end{enumerate}

\vspace*{10pt}

\begin{tcolorbox}[colback=white, colframe=scarlet, title={\textbf{Solution}}, breakable, grow to left by=6mm, grow to right by=10mm]
\begin{enumerate}[label=\textbf{(\alph*)}]
  \item We prove the formula by mathematical induction on $n$.

  \textbf{Base case: $n=1$}

  \[
  f(x)=(1+x)^{-1/2}
  \]

  Differentiating once,

  \[
  f'(x)=-\frac12(1+x)^{-3/2}
  \]

  The given formula gives, when $n=1$,

  \[
  \frac{(-1)^1(2)!}{2^{2}\cdot 1!\,(1+x)^{1+\frac12}}
  =\frac{-2}{4(1+x)^{3/2}}
  =-\frac12(1+x)^{-3/2}
  \]

  So the result is true for $n=1$.

  \textbf{Induction hypothesis}

  Assume that for some $k\ge 1$,

  \[
  f^{(k)}(x)=\frac{(-1)^k(2k)!}{2^{2k}k!(1+x)^{k+\frac12}}
  \]

  \textbf{Induction step}

  Differentiate both sides:

  \begin{align*}
  f^{(k+1)}(x)
  &=\frac{(-1)^k(2k)!}{2^{2k}k!}\cdot \frac{\mathrm{d}}{\mathrm{d}x}\left((1+x)^{-k-\frac12}\right) \\
  &=\frac{(-1)^k(2k)!}{2^{2k}k!}\left(-k-\frac12\right)(1+x)^{-k-\frac32} \\
  &=\frac{(-1)^{k+1}(2k)!}{2^{2k}k!}\left(k+\frac12\right)(1+x)^{-k-\frac32} \\
  &=\frac{(-1)^{k+1}(2k)!}{2^{2k}k!}\cdot \frac{2k+1}{2}(1+x)^{-k-\frac32} \\
  &=\frac{(-1)^{k+1}(2k)!(2k+1)}{2^{2k+1}k!}(1+x)^{-k-\frac32}
  \end{align*}

  Now simplify the coefficient:

  \begin{align*}
  \frac{(2k)!(2k+1)}{2^{2k+1}k!}
  &=\frac{2(k+1)(2k+1)(2k)!}{2^{2k+2}(k+1)k!} \\
  &=\frac{(2k+2)(2k+1)(2k)!}{2^{2k+2}(k+1)!} \\
  &=\frac{(2k+2)!}{2^{2k+2}(k+1)!}
  \end{align*}

  Hence

  \[
  f^{(k+1)}(x)=\frac{(-1)^{k+1}(2k+2)!}{2^{2k+2}(k+1)!(1+x)^{k+\frac32}}
  \]

  which is exactly the required formula with $n=k+1$.

  Therefore, by mathematical induction,

  \[
  f^{(n)}(x)=\frac{(-1)^n(2n)!}{2^{2n}n!(1+x)^{n+\frac12}}
  \qquad \text{for all } n\ge 1
  \]

  \item Using the Maclaurin series,

  \[
  f(x)=f(0)+\sum_{n=1}^{\infty}\frac{f^{(n)}(0)}{n!}x^n
  \]

  First,

  \[
  f(0)=(1+0)^{-1/2}=1
  \]

  From part (a), for $n\ge1$,

  \[
  f^{(n)}(0)=\frac{(-1)^n(2n)!}{2^{2n}n!}
  \]

  So

  \begin{align*}
  (1+x)^{-1/2}
  &=1+\sum_{n=1}^{\infty}\frac{1}{n!}\cdot \frac{(-1)^n(2n)!}{2^{2n}n!}x^n \\
  &=1+\sum_{n=1}^{\infty}(-1)^n\frac{(2n)!}{2^{2n}(n!)^2}x^n
  \end{align*}

  Hence

  \[
  (1+x)^{-1/2}
  =1-\frac{x}{2}+\frac{3x^2}{8}-\frac{5x^3}{16}+\cdots+(-1)^n\frac{(2n)!}{2^{2n}(n!)^2}x^n+\cdots
  \]

\end{enumerate}
\end{tcolorbox}


\newpage

% ---- Question 7 ----
\questionitem

It is given that \[y = (ax + 1)^2 \ln(ax + 1)\] where $a$ is a positive constant. Prove by mathematical induction that, for every integer $n \ge 3$,
\[
\frac{\mathrm{d}^n y}{\mathrm{d}x^n} = 2(-1)^{n-3} \frac{(n-3)!a^n}{(ax + 1)^{n-2}}
\] \marks[-20pt]{6}

\vspace*{10pt}

\begin{tcolorbox}[colback=white, colframe=scarlet, title={\textbf{Solution}}, breakable, grow to left by=6mm, grow to right by=10mm]
Let
\[
P(n): \quad \frac{\mathrm{d}^n y}{\mathrm{d}x^n}=2(-1)^{n-3}\frac{(n-3)!a^n}{(ax+1)^{\,n-2}}
\]
We prove \(P(n)\) is true for all integers \(n \ge 3\) by mathematical induction.

\textbf{Base case: \(n=3\).}

Given
\[
y=(ax+1)^2\ln(ax+1)
\]

Differentiate step by step.

\begin{align*}
\frac{\mathrm{d}y}{\mathrm{d}x}
&=\frac{\mathrm{d}}{\mathrm{d}x}\bigl((ax+1)^2\bigr)\ln(ax+1)+(ax+1)^2\frac{\mathrm{d}}{\mathrm{d}x}\bigl(\ln(ax+1)\bigr)\\
&=2(ax+1)\cdot a\cdot \ln(ax+1)+(ax+1)^2\cdot \frac{a}{ax+1}\\
&=2a(ax+1)\ln(ax+1)+a(ax+1)\\
&=a(ax+1)\bigl(2\ln(ax+1)+1\bigr)
\end{align*}

Now differentiate again.

\begin{align*}
\frac{\mathrm{d}^2y}{\mathrm{d}x^2}
&=a\frac{\mathrm{d}}{\mathrm{d}x}\left[(ax+1)\bigl(2\ln(ax+1)+1\bigr)\right]\\
&=a\left[\frac{\mathrm{d}}{\mathrm{d}x}(ax+1)\bigl(2\ln(ax+1)+1\bigr)+(ax+1)\frac{\mathrm{d}}{\mathrm{d}x}\bigl(2\ln(ax+1)+1\bigr)\right]\\
&=a\left[a\bigl(2\ln(ax+1)+1\bigr)+(ax+1)\left(2\cdot \frac{a}{ax+1}\right)\right]\\
&=a\left[a\bigl(2\ln(ax+1)+1\bigr)+2a\right]\\
&=a^2\bigl(2\ln(ax+1)+3\bigr)
\end{align*}

Differentiate a third time.

\begin{align*}
\frac{\mathrm{d}^3y}{\mathrm{d}x^3}
&=a^2\frac{\mathrm{d}}{\mathrm{d}x}\bigl(2\ln(ax+1)+3\bigr)\\
&=a^2\left(2\cdot \frac{a}{ax+1}\right)\\
&=\frac{2a^3}{ax+1}
\end{align*}

Now substitute \(n=3\) into the claimed formula.

\begin{align*}
2(-1)^{3-3}\frac{(3-3)!a^3}{(ax+1)^{3-2}}
&=2(-1)^0\frac{0!\,a^3}{(ax+1)^1}\\
&=\frac{2a^3}{ax+1}
\end{align*}

This agrees with \(\dfrac{\mathrm{d}^3y}{\mathrm{d}x^3}\), so \(P(3)\) is true.

\textbf{Induction hypothesis.}

Assume that for some integer \(k\ge 3\),
\[
\frac{\mathrm{d}^k y}{\mathrm{d}x^k}=2(-1)^{k-3}\frac{(k-3)!a^k}{(ax+1)^{\,k-2}}
\]

\textbf{Inductive step.}

We must show that this implies
\[
\frac{\mathrm{d}^{k+1} y}{\mathrm{d}x^{k+1}}=2(-1)^{k-2}\frac{(k-2)!a^{k+1}}{(ax+1)^{\,k-1}}
\]

Starting from the induction hypothesis,

\begin{align*}
\frac{\mathrm{d}^{k+1} y}{\mathrm{d}x^{k+1}}
&=\frac{\mathrm{d}}{\mathrm{d}x}\left(2(-1)^{k-3}(k-3)!a^k(ax+1)^{-(k-2)}\right)\\
&=2(-1)^{k-3}(k-3)!a^k\frac{\mathrm{d}}{\mathrm{d}x}\left((ax+1)^{-(k-2)}\right)\\
&=2(-1)^{k-3}(k-3)!a^k\left(-(k-2)(ax+1)^{-(k-1)}\cdot a\right)\\
&=2\bigl[-(-1)^{k-3}\bigr](k-2)(k-3)!a^{k+1}(ax+1)^{-(k-1)}\\
&=2(-1)^{k-2}(k-2)!a^{k+1}(ax+1)^{-(k-1)}\\
&=2(-1)^{k-2}\frac{(k-2)!a^{k+1}}{(ax+1)^{\,k-1}}
\end{align*}

This is exactly the required form for \(P(k+1)\).

So, assuming \(P(k)\) is true has allowed us to prove \(P(k+1)\) is true.

\textbf{Conclusion.}

Since \(P(3)\) is true, and \(P(k)\Rightarrow P(k+1)\) for all integers \(k\ge 3\), it follows by mathematical induction that
\[
\frac{\mathrm{d}^n y}{\mathrm{d}x^n}=2(-1)^{n-3}\frac{(n-3)!a^n}{(ax+1)^{\,n-2}}
\qquad \text{for all integers } n\ge 3
\]

Hence the result is proved.
\end{tcolorbox}


\newpage

% ---- Question 8 ----
\questionitem

The function $f$ is such that $f''(x)=-f(x)$. \\

Prove by mathematical induction that, for every positive integer $n$,
\[
\frac{\mathrm{d}^{2n}}{\mathrm{d}x^{2n}}\left(x^2f(x)\right)=(-1)^n\left(x^2f(x)-4nxf'(x)-2n(2n-1)f(x)\right)
\] \marks[-30pt]{7}

\vspace*{10pt}

\begin{tcolorbox}[colback=white, colframe=scarlet, title={\textbf{Solution}}, breakable, grow to left by=6mm, grow to right by=10mm]
We prove the result by mathematical induction on \(n\).

Let \(P(n)\) be the statement
\[
\frac{\mathrm{d}^{2n}}{\mathrm{d}x^{2n}}\left(x^2f(x)\right)=(-1)^n\left(x^2f(x)-4nxf'(x)-2n(2n-1)f(x)\right)
\]

\textbf{Base case: \(n=1\)}

We calculate the second derivative directly:
\begin{align*}
\frac{\mathrm{d}}{\mathrm{d}x}\left(x^2f(x)\right)
&=2xf(x)+x^2f'(x)
\\
\frac{\mathrm{d}^2}{\mathrm{d}x^2}\left(x^2f(x)\right)
&=\frac{\mathrm{d}}{\mathrm{d}x}\left(2xf+x^2f'\right)
\\
&=2f+2xf'+2xf'+x^2f''
\\
&=2f+4xf'+x^2f''
\end{align*}

Given that \(f''(x)=-f(x)\), this becomes
\begin{align*}
\frac{\mathrm{d}^2}{\mathrm{d}x^2}\left(x^2f(x)\right)
&=2f+4xf'-x^2f
\\
&=-(x^2f-4xf'-2f)
\\
&=(-1)^1\left(x^2f-4(1)xf'-2(1)(2\cdot 1-1)f\right)
\end{align*}

So \(P(1)\) is true.

\textbf{Inductive hypothesis}

Assume that \(P(k)\) is true for some positive integer \(k\). Then
\[
\frac{\mathrm{d}^{2k}}{\mathrm{d}x^{2k}}\left(x^2f\right)
= (-1)^k\left(x^2f-4kxf'-2k(2k-1)f\right)
\]

\textbf{Inductive step}

We must show that \(P(k+1)\) is true.

Starting from the inductive hypothesis and differentiating twice,
\begin{align*}
\frac{\mathrm{d}^{2k+2}}{\mathrm{d}x^{2k+2}}\left(x^2f\right)
&= \frac{\mathrm{d}^2}{\mathrm{d}x^2}\left[\frac{\mathrm{d}^{2k}}{\mathrm{d}x^{2k}}(x^2f)\right]
\\
&= (-1)^k\frac{\mathrm{d}^2}{\mathrm{d}x^2}\left(x^2f-4kxf'-2k(2k-1)f\right)
\\
&= (-1)^k\left((x^2f)''-4k(xf')''-2k(2k-1)f''\right)
\end{align*}

Now work out each of these second derivatives.

First, from the base case calculation,
\[
(x^2f)''=-x^2f+4xf'+2f
\]

Next,
\begin{align*}
(xf')'
&=f'+xf''
\\
(xf')''
&=f''+f''+xf'''
\\
&=2f''+xf'''
\end{align*}

Since \(f''=-f\), differentiating gives \(f'''=-f'\). Therefore
\[
(xf')''=2(-f)+x(-f')=-2f-xf'
\]
so
\[
(xf')''=-xf'-2f
\]

Also,
\[
f''=-f
\]

Substituting these into the inductive step,
\begin{align*}
\frac{\mathrm{d}^{2k+2}}{\mathrm{d}x^{2k+2}}\left(x^2f\right)
&= (-1)^k\left[(-x^2f+4xf'+2f)-4k(-xf'-2f)-2k(2k-1)(-f)\right]
\\
&= (-1)^k\left[-x^2f+4xf'+2f+4kxf'+8kf+2k(2k-1)f\right]
\\
&= (-1)^k\left[-x^2f+4(k+1)xf'+\bigl(2+8k+2k(2k-1)\bigr)f\right]
\end{align*}

Now simplify the coefficient of \(f\):
\begin{align*}
2+8k+2k(2k-1)
&=2+8k+4k^2-2k
\\
&=4k^2+6k+2
\\
&=2(k+1)(2k+1)
\end{align*}

Hence
\begin{align*}
\frac{\mathrm{d}^{2k+2}}{\mathrm{d}x^{2k+2}}\left(x^2f\right)
&= (-1)^k\left[-x^2f+4(k+1)xf'+2(k+1)(2k+1)f\right]
\\
&= (-1)^{k+1}\left[x^2f-4(k+1)xf'-2(k+1)(2k+1)f\right]
\end{align*}

Since \(2(k+1)-1=2k+1\), this is exactly
\[
\frac{\mathrm{d}^{2(k+1)}}{\mathrm{d}x^{2(k+1)}}\left(x^2f\right)
=
(-1)^{k+1}\left(x^2f-4(k+1)xf'-2(k+1)(2(k+1)-1)f\right)
\]

So \(P(k+1)\) is true.

\textbf{Conclusion}

\(P(1)\) is true, and \(P(k)\Rightarrow P(k+1)\). Therefore, by mathematical induction,
\[
\frac{\mathrm{d}^{2n}}{\mathrm{d}x^{2n}}\left(x^2f(x)\right)=(-1)^n\left(x^2f(x)-4nxf'(x)-2n(2n-1)f(x)\right)
\]
for every positive integer \(n\).
\end{tcolorbox}


\end{enumerate}
\end{document}
